\section{Open Questions}

\begin{enumerate}
  \item \textbf{Segmented BK / parity mapping variants.} The paper analyses only
        standard BK and JW. Does the segmented BK (BKSF) variant reproduce the
        same H$_2$ gate counts, and where does its locality/CNOT tradeoff cross
        over versus standard BK as $n$ grows?
        \emph{Next step:} reimplement BKSF and parity encodings using the same
        recursive-matrix scaffolding in \texttt{replicate.py}; assemble the H$_2$
        Pauli Hamiltonian in each; compute Trotter-step gate counts for
        $n\in\{4,8,16,32,64\}$ and plot BK, JW, BKSF, parity side by side.

  \item \textbf{Qubit tapering via $Z_2$ symmetries.} Tapering was not yet
        standard when the paper appeared; how much does modern
        particle-number/spin/point-group tapering reduce the BK H$_2$ gate
        counts below the 30 sq / 44 CNOT baseline reported here?
        \emph{Next step:} apply an OpenFermion or hand-rolled $Z_2$ taper to the
        BK H$_2$ Hamiltonian, recount sq and CNOT for one Trotter step,
        and compare.

  \item \textbf{Noise robustness of BK-encoded operators.} The paper's gate-count
        comparison is noiseless. Does the BK $O(\log n)$ locality advantage
        translate into a measurable NISQ-era simulation-fidelity advantage on a
        realistic noise model, or is the constant-factor CNOT overhead
        dominant?
        \emph{Next step:} run \texttt{qiskit-aer} noisy simulations of one BK vs.\
        JW Trotter step at several depolarising + T1/T2 noise levels calibrated
        to a public IBM backend snapshot; plot fidelity vs.\ two-qubit error rate.

  \item \textbf{Scaling to larger chemistry benchmarks.} For LiH, BeH$_2$, H$_2$O
        in STO-3G/6-31G, does the observed H$_2$ BK/JW ratio generalise, or does
        BK's log-locality compound to a total-gate-cost win in larger active
        spaces?
        \emph{Next step:} generate Pauli Hamiltonians in both encodings via
        PySCF+OpenFermion, compile Trotter steps with the textbook rule used
        here, and tabulate gate counts as active-space size grows from 4 to
        about 20 spin-orbitals.

  \item \textbf{Comparison with post-2012 fermion mappings.} How does BK compare
        to compact encodings, Verstraete--Cirac, ternary-tree, and Derby--Klassen
        on the same H$_2$ benchmark used in the paper?
        \emph{Next step:} reimplement or import each mapping, build H$_2$ Pauli
        Hamiltonians, and compare qubit count, string count, per-operator
        locality, and one-Trotter-step sq/CNOT counts in a single table.
\end{enumerate}
